%%=============================================================================
%% Samenvatting
%%=============================================================================

% TODO: De "abstract" of samenvatting is een kernachtige (~ 1 blz. voor een
% thesis) synthese van het document.
%
% Een goede abstract biedt een kernachtig antwoord op volgende vragen:
%
% 1. Waarover gaat de bachelorproef?
% 2. Waarom heb je er over geschreven?
% 3. Hoe heb je het onderzoek uitgevoerd?
% 4. Wat waren de resultaten? Wat blijkt uit je onderzoek?
% 5. Wat betekenen je resultaten? Wat is de relevantie voor het werkveld?
%
% Daarom bestaat een abstract uit volgende componenten:
%
% - inleiding + kaderen thema
% - probleemstelling
% - (centrale) onderzoeksvraag
% - onderzoeksdoelstelling
% - methodologie
% - resultaten (beperk tot de belangrijkste, relevant voor de onderzoeksvraag)
% - conclusies, aanbevelingen, beperkingen
%
% LET OP! Een samenvatting is GEEN voorwoord!

%%---------- Nederlandse samenvatting -----------------------------------------
%
% TODO: Als je je bachelorproef in het Engels schrijft, moet je eerst een
% Nederlandse samenvatting invoegen. Haal daarvoor onderstaande code uit
% commentaar.
% Wie zijn bachelorproef in het Nederlands schrijft, kan dit negeren, de inhoud
% wordt niet in het document ingevoegd.

\IfLanguageName{english}{%
\selectlanguage{dutch}
\chapter*{Samenvatting}
\lipsum[1-4]
\selectlanguage{english}
}{}

%%---------- Samenvatting -----------------------------------------------------
% De samenvatting in de hoofdtaal van het document

\chapter*{\IfLanguageName{dutch}{Samenvatting}{Abstract}}

Deze bachelorproef behandelt het ontwerp en de realisatie van RustProv, een provisioner voor virtuele Windows- en Linuxomgevingen. Binnen de opleiding Toegepaste Informatica aan HOGENT wordt voor het opzetten van labo- en ontwikkelomgevingen vaak gebruikgemaakt van Vagrant, maar in de praktijk leidt dit soms tot trage opstarttijden, foutgevoelige provisioning en beperkte controle over het proces. Vanu    it dit punt werd onderzocht hoe een eigen provisioner in Rust ontwikkeld kan worden die beter aansluit bij de benodigdheden van een onderwijsomgeving.

Het doel van deze bachelorproef is een proof of concept te realiseren dat virtuele machines automatisch kan aanmaken, configureren en provisionen op een consistente en betrouwbare manier. Daarbij wordt ook gekeken naar praktische verbeteringen, zoals het direct koppelen van een VDI-bestand, het gebruik van shared folders en het ondersteunen van zowel één VM als meerdere VM's. De performantie van RustProv wordt bovendien vergeleken met die van Vagrant om na te gaan of de nieuwe aanpak sneller of efficiënter is.

De experimenten tonen aan dat RustProv in de meeste scenario's sneller is dan Vagrant. Bij drie Linux-VM's is RustProv gemiddeld 74,03 seconden sneller en bij drie Windows-VM's bedraagt de tijdswinst gemiddeld 472,10 seconden. Bij één Linux-VM is Vagrant iets sneller, maar bij één Windows-VM is RustProv 58,91 seconden sneller. Deze resultaten bevestigen dat een lichtgewicht, in Rust ontwikkelde provisioner concurrentievoordelen kan bieden ten opzichte van Vagrant, vooral bij het provisionen van meerdere VM's.

Deze bachelorproef levert daarmee een werkende basis op voor verdere uitbreiding van geautomatiseerde VM-provisioning in een onderwijscontext.
